\section{代数系统}

\subsection{代数结构}

\subsubsection{基本概念}

\begin{enumerate}
	\item \textbf{代数系统}：由非空集合及其上的若干运算组成，记作
	\[
	\langle A,f_1,f_2,\dots,f_k\rangle.
	\]
	\item \textbf{\(n\) 元运算}：从 \(A^n\) 到 \(A\) 的映射称为 \(A\) 上的 \(n\) 元运算。二元运算常写作 \(a*b\)。
	\item \textbf{封闭性}：若
	\[
	\forall a,b\in A,\quad a*b\in A,
	\]
	则称 \(*\) 在 \(A\) 上封闭。
	\item \textbf{运算表}：有限代数系统可用 Cayley 表分析。封闭性对应表中结果全在集合内；交换律对应表关于主对角线对称。
\end{enumerate}

\subsubsection{二元运算的常见性质}

设 \(*,\circ\) 是集合 \(A\) 上的二元运算。
\begin{enumerate}
	\item 交换律：
	\[
	a*b=b*a.
	\]
	\item 结合律：
	\[
	(a*b)*c=a*(b*c).
	\]
	\item 分配律：
	\[
	a*(b\circ c)=(a*b)\circ(a*c),
	\]
	\[
	(b\circ c)*a=(b*a)\circ(c*a).
	\]
	\item 吸收律：
	\[
	a*(a\circ b)=a,\qquad a\circ(a*b)=a.
	\]
	\item 幂等律：
	\[
	a*a=a.
	\]
\end{enumerate}

\subsubsection{特殊元素}

\begin{enumerate}
	\item \textbf{零元}：若存在 \(\theta\in A\)，使
	\[
	a*\theta=\theta*a=\theta,\qquad \forall a\in A,
	\]
	则 \(\theta\) 为零元。
	\item \textbf{幺元}：若存在 \(e\in A\)，使
	\[
	a*e=e*a=a,\qquad \forall a\in A,
	\]
	则 \(e\) 为幺元或单位元。
	\item \textbf{逆元}：若 \(e\) 是幺元，且
	\[
	a*b=b*a=e,
	\]
	则称 \(b\) 是 \(a\) 的逆元，记作 \(a^{-1}\)。
	\item 左、右零元都存在时二者相等；左、右幺元都存在时二者相等。
\end{enumerate}

\subsection{半群、群、子群}

\subsubsection{半群与独异点}

\begin{enumerate}
	\item \textbf{广群}：非空集合 \(G\) 上的二元运算 \(*\) 封闭，则 \(\langle G,*\rangle\) 称为广群。
	\item \textbf{半群}：满足封闭性和结合律的代数系统称为半群。
	\item \textbf{独异点}：含有幺元的半群称为独异点，也称含幺半群。
\end{enumerate}

\subsubsection{群}

若 \(\langle G,*\rangle\) 是独异点，且每个元素均有逆元，则称其为群。若群的运算满足交换律，则称为阿贝尔群或交换群。

群中的常用结论：
\begin{enumerate}
	\item 幺元唯一。
	\item 每个元素的逆元唯一。
	\item \((a^{-1})^{-1}=a\)。
	\item \((a*b)^{-1}=b^{-1}*a^{-1}\)。
	\item 消去律：
	\[
	a*b=a*c\Rightarrow b=c,\qquad b*a=c*a\Rightarrow b=c.
	\]
	\item 群中除幺元外没有其他幂等元。若 \(a*a=a\)，则 \(a=e\)。
	\item 有限群的运算表中，每一行、每一列都是群元素的一个排列。
\end{enumerate}

\begin{question}{例题1}
	判断 \(\langle\mathbb{Z},+\rangle\) 是否为群。
	
	\begin{solution}
	整数加法封闭并满足结合律；幺元为 \(0\)；任意 \(a\in\mathbb{Z}\) 的逆元为 \(-a\)。因此 \(\langle\mathbb{Z},+\rangle\) 是群。又因为加法满足交换律，所以它是阿贝尔群。
	\end{solution}
\end{question}

\subsubsection{子群}

若 \(H\subseteq G\)，且 \(\langle H,*\rangle\) 本身也是群，则称 \(H\) 是 \(G\) 的子群，记作 \(H\le G\)。

子群判定常用结论：
\begin{enumerate}
	\item 非空子集 \(H\subseteq G\) 是子群，当且仅当
	\[
	\forall a,b\in H,\quad a*b^{-1}\in H.
	\]
	\item 若 \(G\) 是有限群，非空子集 \(H\subseteq G\) 对运算封闭，则 \(H\le G\)。
\end{enumerate}

\subsection{阿贝尔群与循环群}

\subsubsection{循环群}

设 \(G\) 是群，若存在 \(a\in G\)，使得
\[
G=\langle a\rangle=\{a^n\mid n\in\mathbb{Z}\},
\]
则称 \(G\) 是循环群，\(a\) 称为生成元。

\begin{enumerate}
	\item 任意循环群都是阿贝尔群。
	\item 阿贝尔群未必是循环群。
	\item 元素 \(a\) 的阶是使 \(a^n=e\) 成立的最小正整数 \(n\)。若不存在这样的正整数，则称 \(a\) 为无限阶。
\end{enumerate}

\subsubsection{置换群}

集合 \(S\) 到自身的双射称为 \(S\) 上的一个置换。所有置换在函数复合下构成群，称为置换群。集合 \(\{1,2,\dots,n\}\) 上的全体置换群记作 \(S_n\)，称为 \(n\) 次对称群。

\subsection{陪集和拉格朗日定理}

设 \(G\) 是群，\(H\le G\)，\(a\in G\)。

\begin{enumerate}
	\item 左陪集：
	\[
	aH=\{a*h\mid h\in H\}.
	\]
	\item 右陪集：
	\[
	Ha=\{h*a\mid h\in H\}.
	\]
	\item \(H\) 在 \(G\) 中不同左陪集的个数称为 \(H\) 在 \(G\) 中的指数，记作 \([G:H]\)。
\end{enumerate}

\textbf{拉格朗日定理}：若 \(G\) 是有限群，\(H\le G\)，则
\[
|G|=[G:H]\cdot |H|.
\]
因此 \(|H|\) 整除 \(|G|\)。

重要推论：
\begin{enumerate}
	\item 有限群中任意元素的阶都整除群的阶。
	\item 若 \(|G|=n\)，则对任意 \(a\in G\)，有 \(a^n=e\)。
	\item 素数阶群一定是循环群，且任一非幺元都是生成元。
\end{enumerate}

\begin{question}{例题2}
	设 \(G\) 是 \(12\) 阶有限群，\(H\le G\)，且 \(|H|=3\)。求 \([G:H]\)。
	
	\begin{solution}
	由拉格朗日定理，
	\[
	[G:H]=\frac{|G|}{|H|}=\frac{12}{3}=4.
	\]
	\end{solution}
\end{question}

\subsection{同态和同构}

\subsubsection{同态}

设 \(\langle A,*\rangle\)、\(\langle B,\circ\rangle\) 是两个代数系统。映射 \(f:A\to B\) 若满足
\[
f(x*y)=f(x)\circ f(y),\qquad \forall x,y\in A,
\]
则称 \(f\) 为同态映射。

\begin{enumerate}
	\item 同态若为单射，称为单同态。
	\item 同态若为满射，称为满同态。
	\item 同态若为双射，称为同构。
	\item 若 \(A\) 与 \(B\) 同构，记作 \(A\cong B\)。同构表示两个代数系统结构相同。
\end{enumerate}

群同态 \(f:G\to G'\) 的核为
\[
\operatorname{Ker}f=\{x\in G\mid f(x)=e'\},
\]
其中 \(e'\) 是 \(G'\) 的幺元。

\subsection{环与域}

\subsubsection{环}

代数系统 \(\langle A,+,\cdot\rangle\) 若满足：
\begin{enumerate}
	\item \(\langle A,+\rangle\) 是阿贝尔群；
	\item \(\langle A,\cdot\rangle\) 是半群；
	\item 乘法对加法满足左右分配律：
	\[
	a\cdot(b+c)=a\cdot b+a\cdot c,
	\]
	\[
	(b+c)\cdot a=b\cdot a+c\cdot a,
	\]
\end{enumerate}
则称 \(\langle A,+,\cdot\rangle\) 为环。

\subsubsection{整环与域}

\begin{enumerate}
	\item \textbf{整环}：乘法可交换、有乘法幺元且无零因子的环称为整环。
	\item \textbf{域}：若 \(\langle A,+\rangle\) 是阿贝尔群，\(A-\{0\}\) 在乘法下是阿贝尔群，且乘法对加法满足分配律，则称为域。
	\item 域一定是整环。
	\item 有限整环一定是域。
\end{enumerate}

\begin{center}
	\begin{tabular}{cc}
		\toprule
		结构 & 典型例子\\
		\midrule
		环 & \(\langle\mathbb{Z},+,\cdot\rangle\)\\
		整环 & \(\langle\mathbb{Z},+,\cdot\rangle\)\\
		域 & \(\langle\mathbb{Q},+,\cdot\rangle,\ \langle\mathbb{R},+,\cdot\rangle\)\\
		非交换环 & \(n\) 阶矩阵环\\
		\bottomrule
	\end{tabular}
\end{center}

\section{格和布尔代数}

\subsection{格的基本概念}

\subsubsection{格}

设 \(\langle A,\le\rangle\) 是偏序集。若任意 \(a,b\in A\) 都有最小上界和最大下界，则称 \(\langle A,\le\rangle\) 为格。

\begin{enumerate}
	\item \(a\vee b\)：\(a,b\) 的最小上界，称为并。
	\item \(a\wedge b\)：\(a,b\) 的最大下界，称为交。
	\item 全序集一定是格，但偏序集未必都是格。
\end{enumerate}

\subsubsection{格的代数性质}

若 \(\langle A,\le\rangle\) 是格，则 \(\langle A,\vee,\wedge\rangle\) 满足：
\begin{enumerate}
	\item 交换律：
	\[
	a\vee b=b\vee a,\qquad a\wedge b=b\wedge a.
	\]
	\item 结合律：
	\[
	a\vee(b\vee c)=(a\vee b)\vee c,
	\]
	\[
	a\wedge(b\wedge c)=(a\wedge b)\wedge c.
	\]
	\item 吸收律：
	\[
	a\vee(a\wedge b)=a,\qquad a\wedge(a\vee b)=a.
	\]
	\item 幂等律：
	\[
	a\vee a=a,\qquad a\wedge a=a.
	\]
\end{enumerate}

反过来，若 \(\langle A,\vee,\wedge\rangle\) 满足交换律、结合律和吸收律，则可定义偏序：
\[
a\le b\Leftrightarrow a\wedge b=a,
\]
也等价于
\[
a\le b\Leftrightarrow a\vee b=b.
\]

\subsubsection{格的对偶原则}

把 \(\vee\) 与 \(\wedge\) 互换，把 \(0\) 与 \(1\) 互换，把 \(\le\) 与 \(\ge\) 互换，所得命题称为原命题的对偶命题。若原命题在所有格中成立，则其对偶命题也成立。

\subsection{特殊的格}

\subsubsection{分配格}

若格满足
\[
a\vee(b\wedge c)=(a\vee b)\wedge(a\vee c),
\]
\[
a\wedge(b\vee c)=(a\wedge b)\vee(a\wedge c),
\]
则称为分配格。链一定是分配格。

任意格中总有分配不等式：
\[
a\vee(b\wedge c)\le (a\vee b)\wedge(a\vee c),
\]
\[
(a\wedge b)\vee(a\wedge c)\le a\wedge(b\vee c).
\]

\subsubsection{有界格、有补格、布尔格}

\begin{enumerate}
	\item \textbf{有界格}：存在全下界 \(0\) 和全上界 \(1\)，使
	\[
	0\le a\le1,\qquad \forall a\in A.
	\]
	\item \textbf{补元}：在有界格中，若
	\[
	a\vee b=1,\qquad a\wedge b=0,
	\]
	则称 \(b\) 是 \(a\) 的补元。
	\item \textbf{有补格}：每个元素都有补元的有界格。
	\item \textbf{布尔格}：有补分配格。
\end{enumerate}

\begin{question}{例题3}
	证明 \(\langle\mathcal{P}(S),\subseteq\rangle\) 是布尔格。
	
	\begin{solution}
	对任意 \(A,B\subseteq S\)，最小上界是 \(A\cup B\)，最大下界是 \(A\cap B\)，因此它是格。全下界为 \(\varnothing\)，全上界为 \(S\)。集合并交满足分配律。任意 \(A\subseteq S\) 的补元为 \(S-A\)，因为
	\[
	A\cup(S-A)=S,\qquad A\cap(S-A)=\varnothing.
	\]
	所以 \(\langle\mathcal{P}(S),\subseteq\rangle\) 是布尔格。
	\end{solution}
\end{question}

\subsection{布尔代数}

\subsubsection{布尔代数定义}

由布尔格诱导的代数系统
\[
\langle A,\vee,\wedge,{}'\!,0,1\rangle
\]
称为布尔代数，其中 \(a'\) 表示 \(a\) 的补元。

\subsubsection{布尔代数常用公式}

\begin{center}
	\begin{tabular}{cl}
		\toprule
		名称 & 公式\\
		\midrule
		同一律 & \(a\vee0=a,\quad a\wedge1=a\)\\
		零律 & \(a\vee1=1,\quad a\wedge0=0\)\\
		补元律 & \(a\vee a'=1,\quad a\wedge a'=0\)\\
		双补律 & \((a')'=a\)\\
		德摩根律 & \((a\vee b)'=a'\wedge b'\)\\
		          & \((a\wedge b)'=a'\vee b'\)\\
		\bottomrule
	\end{tabular}
\end{center}

\subsubsection{有限布尔代数与原子}

\begin{enumerate}
	\item 若 \(a\ne0\)，且不存在 \(x\) 满足 \(0<x<a\)，则称 \(a\) 为原子。
	\item 在有限布尔代数中，每个非零元素都可表示为若干原子的并。
	\item 若有限布尔代数有 \(n\) 个原子，则它共有 \(2^n\) 个元素。
\end{enumerate}

\begin{question}{例题4}
	设 \(S=\{1,2,3\}\)。布尔代数 \(\mathcal{P}(S)\) 有几个原子？共有几个元素？
	
	\begin{solution}
	原子为单元素子集
	\[
	\{1\},\{2\},\{3\}.
	\]
	所以原子数为 \(3\)，元素数为 \(2^3=8\)。
	\end{solution}
\end{question}

\section{代数系统考试补充}

\subsection{运算表判定技巧}

有限代数系统的运算表是考试高频题型。设运算表行表示左操作数，列表示右操作数。
\begin{enumerate}
	\item \textbf{封闭性}：表中每个结果都必须属于给定集合。
	\item \textbf{交换律}：运算表关于主对角线对称。
	\item \textbf{幺元}：若某一行与表头完全相同，且同名列也与表头完全相同，则该元素为幺元。
	\item \textbf{零元}：若某一行全为该元素，且同名列也全为该元素，则该元素为零元。
	\item \textbf{逆元}：在有幺元 \(e\) 的前提下，若 \(a*b=b*a=e\)，则 \(a,b\) 互为逆元。
	\item \textbf{群表}：有限群的每一行、每一列都必须是集合元素的一个排列。
\end{enumerate}

\begin{question}{例题5}
	设 \(A=\{a,b,c\}\)，运算 \(*\) 的表为
	\[
	\begin{array}{cccc}
	\toprule
	*&a&b&c\\
	\midrule
	a&a&b&c\\
	b&b&a&c\\
	c&c&c&c\\
	\bottomrule
	\end{array}
	\]
	求幺元和零元。
	
	\begin{solution}
	第 \(a\) 行为 \(a,b,c\)，第 \(a\) 列也为 \(a,b,c\)，所以 \(a\) 是幺元。第 \(c\) 行全为 \(c\)，第 \(c\) 列也全为 \(c\)，所以 \(c\) 是零元。
	\end{solution}
\end{question}

\subsection{模 n 运算}

整数模 \(n\) 同余：
\[
a\equiv b\pmod n\Leftrightarrow n\mid(a-b).
\]
同余模 \(n\) 是 \(\mathbb Z\) 上的等价关系，其等价类为
\[
[0],[1],\dots,[n-1].
\]

集合 \(\mathbb Z_n=\{[0],[1],\dots,[n-1]\}\) 上可定义模 \(n\) 加法和乘法：
\[
[a]+[b]=[a+b],\qquad [a]\cdot[b]=[ab].
\]

\begin{enumerate}
	\item \(\langle\mathbb Z_n,+\rangle\) 是 \(n\) 阶循环群，生成元为 \([1]\)。
	\item \(\langle\mathbb Z_n,+,\cdot\rangle\) 是有幺交换环。
	\item 当且仅当 \(p\) 是素数时，\(\mathbb Z_p\) 是域。
\end{enumerate}

\subsection{布尔函数与析取范式}

布尔代数在计算机专业中常用于开关电路与逻辑门。
\begin{enumerate}
	\item 布尔函数是从 \(\{0,1\}^n\) 到 \(\{0,1\}\) 的函数。
	\item \(n\) 元布尔函数共有
	\[
	2^{2^n}
	\]
	个。
	\item 每个布尔函数都可以表示为主析取范式或主合取范式。
\end{enumerate}

若布尔函数在小项 \(m_{i_1},m_{i_2},\dots,m_{i_k}\) 上取值为 \(1\)，则可写作
\[
f=\sum m(i_1,i_2,\dots,i_k).
\]
若在大项 \(M_{j_1},M_{j_2},\dots,M_{j_s}\) 上取值为 \(0\)，则可写作
\[
f=\prod M(j_1,j_2,\dots,j_s).
\]

\subsection{本章重点定义与定理速查}

\subsubsection{代数系统与运算}

\begin{enumerate}
	\item \textbf{代数系统}：非空集合连同其上的若干运算组成的系统，常记为 \(\langle A,*\rangle\) 或 \(\langle A,+,\cdot\rangle\)。
	\item \textbf{封闭性}：对任意 \(a,b\in A\)，都有 \(a*b\in A\)。
	\item \textbf{结合律}：
	\[
	(a*b)*c=a*(b*c).
	\]
	\item \textbf{交换律}：
	\[
	a*b=b*a.
	\]
	\item \textbf{幺元}：若 \(a*e=e*a=a\)，则 \(e\) 是幺元。若存在，幺元唯一。
	\item \textbf{零元}：若 \(a*\theta=\theta*a=\theta\)，则 \(\theta\) 是零元。若存在，零元唯一。
	\item \textbf{逆元}：在有幺元 \(e\) 的系统中，若 \(a*b=b*a=e\)，则 \(b\) 是 \(a\) 的逆元。群中每个元素逆元唯一。
\end{enumerate}

\subsubsection{群论核心定义与定理}

\begin{enumerate}
	\item \textbf{半群}：封闭且满足结合律的代数系统。
	\item \textbf{独异点}：含幺元的半群。
	\item \textbf{群}：含幺元且每个元素都有逆元的半群。
	\item \textbf{阿贝尔群}：运算满足交换律的群。
	\item \textbf{子群判定定理}：非空子集 \(H\subseteq G\) 是子群，当且仅当
	\[
	\forall a,b\in H,\quad ab^{-1}\in H.
	\]
	若 \(G\) 有限，非空子集 \(H\) 对运算封闭即可推出 \(H\le G\)。
	\item \textbf{循环群}：若存在 \(a\in G\)，使
	\[
	G=\langle a\rangle=\{a^n\mid n\in\mathbb Z\},
	\]
	则 \(G\) 是循环群。循环群一定是阿贝尔群。
	\item \textbf{元素的阶}：使 \(a^n=e\) 成立的最小正整数 \(n\)。若不存在，则为无限阶。
	\item \textbf{拉格朗日定理}：若 \(G\) 是有限群，\(H\le G\)，则
	\[
	|G|=[G:H]|H|.
	\]
	因此子群阶整除群阶。
	\item \textbf{拉格朗日推论}：有限群中任意元素的阶整除群的阶；若 \(|G|=n\)，则 \(a^n=e\)。
	\item \textbf{素数阶群定理}：素数阶群一定是循环群，且任意非幺元都是生成元。
\end{enumerate}

\subsubsection{同态、环与域}

\begin{enumerate}
	\item \textbf{同态}：映射 \(f:A\to B\) 保持运算，即
	\[
	f(x*y)=f(x)\circ f(y).
	\]
	\item \textbf{同构}：双射同态。两个代数系统同构表示运算结构完全相同。
	\item \textbf{群同态核}：
	\[
	\operatorname{Ker}f=\{x\in G\mid f(x)=e'\}.
	\]
	核是 \(G\) 的子群。
	\item \textbf{环}：\(\langle R,+\rangle\) 是阿贝尔群，\(\langle R,\cdot\rangle\) 是半群，且乘法对加法满足左右分配律。
	\item \textbf{整环}：有乘法幺元、乘法交换且无零因子的环。
	\item \textbf{域}：非零元素对乘法构成阿贝尔群的交换环。域一定是整环，有限整环一定是域。
	\item \textbf{模 \(n\) 结论}：\(\mathbb Z_n\) 在模加法下是循环群；\(\mathbb Z_n\) 是有幺交换环；\(\mathbb Z_p\) 是域当且仅当 \(p\) 是素数。
\end{enumerate}

\subsubsection{格与布尔代数}

\begin{enumerate}
	\item \textbf{格}：偏序集中任意两个元素都有最小上界和最大下界。
	\item \textbf{并与交}：\(a\vee b\) 表示最小上界，\(a\wedge b\) 表示最大下界。
	\item \textbf{格的代数刻画}：满足交换律、结合律和吸收律的 \(\langle A,\vee,\wedge\rangle\) 可诱导出格序。
	\item \textbf{分配格}：满足
	\[
	a\vee(b\wedge c)=(a\vee b)\wedge(a\vee c),
	\]
	\[
	a\wedge(b\vee c)=(a\wedge b)\vee(a\wedge c).
	\]
	\item \textbf{有界格}：有全下界 \(0\) 和全上界 \(1\) 的格。
	\item \textbf{补元}：若 \(a\vee b=1\) 且 \(a\wedge b=0\)，则 \(b\) 是 \(a\) 的补元。
	\item \textbf{布尔格}：有补分配格。
	\item \textbf{布尔代数}：由布尔格诱导的代数系统 \(\langle A,\vee,\wedge,{}'\!,0,1\rangle\)。
	\item \textbf{有限布尔代数原子定理}：若有限布尔代数有 \(n\) 个原子，则共有 \(2^n\) 个元素。
	\item \textbf{布尔函数计数}：\(n\) 元布尔函数共有
	\[
	2^{2^n}
	\]
	个。
\end{enumerate}
